\section{Structural-Codomain Extension of the Trace-to-Scheme Export Theorem}
\label{sec:structural-codomain-extension}

The export theorem of \S\ref{sec:trace-to-scheme-export-v35} states the six necessary gates G1--G6 for an APF trace value to be physically exported. Its empirical training set (the W on-shell route, the four fermion-sector routes, the charged-lepton residual closeout) is uniformly numerical: each route's codomain is a single number with a propagated uncertainty. The framework's \textsf{closed} Reconstruction Programs include routes of a different kind --- structural-codomain routes whose target is a theorem-shape rather than a number-shape. The Riemannian-geometry closure (RP1) is the canonical example: the codomain is GR's geometric structure, derived from $A_1$ plus algebraic data without any numerical residual to propagate.

This section extends the export theorem to cover structural-codomain routes by giving G4 and G6 their natural structural-form readings. G1, G2, G3, and G5 read identically across numerical and structural forms.

\subsection{Numerical and structural routes}
\label{subsec:numerical-vs-structural}

A route $\mathcal{R}$ is \emph{numerical} if its codomain is a real number $M_{\mathcal{S}} \in \mathbb{R}$ together with a measurement uncertainty $\sigma_M > 0$. A route is \emph{structural} if its codomain is a theorem-shape: an assertion about the structural form of an object derived from APF, holding inside a declared domain of validity and breaking down at named edge cases.

Mass routes are numerical: $M_W^{\rm OS}$, $m_b(m_b)_{\rm MSbar}$, etc. Cosmological budget predictions are numerical: $\eta_B$, $n_s$, $\Omega_b/\Omega_c/\Omega_\Lambda$. Riemannian geometry's emergence is structural: the codomain is ``causal order $\Rightarrow$ conformal class $[g]$'' (HKM 1976) or ``$\delta Q = T\, dS \Rightarrow$ linearized Einstein equations.'' No single number is the codomain.

\subsection{Structural form of G4 (covariance $\to$ domain of validity)}
\label{subsec:G4-structural}

Numerical-form G4 requires the route to propagate input uncertainty through the transport $\tau$ to produce a covariance on $\tau(T)$. For a structural-codomain route there is no numerical $\tau(T)$ to attach a covariance to; the closure is an assertion of structural form, valid uniformly within its domain.

Structural-form G4 reads: \emph{the domain of validity of the structural closure is declared}. The conditions under which the derivation holds are named; the regimes in which it does not extend are explicit. For RP1, the domain of validity is the finite-physical regime within causally connected regions; the closure does not extend to the quantum-gravity regime (parked at RP5 long-term capstone) or to far-from-equilibrium spacetimes (named in Paper 25 H4).

The structural form preserves G4's underlying necessity: $A_1$ requires every distinction's cost to be accounted for. For numerical routes that accounting takes the form of propagated covariance; for structural routes it takes the form of declared domain of validity --- distinctions outside the domain remain unaccounted for and are \emph{by declaration} not part of the route's claim.

\subsection{Structural form of G6 (residual channel $\to$ named edge cases)}
\label{subsec:G6-structural}

Numerical-form G6 requires the post-comparison residual $\rho = \tau(T) - M_{\mathcal{S}}$ to be assigned to a named channel (covariance-consistent, named operator-of-this-route, named scheme convention, named theory-uncertainty band). For a structural-codomain route there is no $\rho$ in the numerical sense; what plays the analogous role is the set of regimes at the boundary of the route's domain of validity.

Structural-form G6 reads: \emph{the edge cases at the derivation's boundary are named}. RP1's structural-form G6 closes by naming the quantum-gravity regime as RP5, the far-from-equilibrium regime as Paper 25 H4, and the non-Lorentzian structural variants as out-of-scope. The framework does not claim its structural derivation extends into these regimes, and it names each regime explicitly so that any reader knows where the closure stops.

The structural form preserves G6's underlying necessity: $A_1 + $ MD require every distinction in the framework's content to have a paying ledger. For numerical routes that means residuals must have a channel; for structural routes it means edge cases must have a parking destination. An unnamed edge case at the derivation's boundary is content the framework has not paid for.

\subsection{Worked example: RP1 spacetime emergence}
\label{subsec:rp1-walk}

We walk RP1's structural closure through G1--G6 in the structural-form interpretation.

\begin{itemize}
\item \textbf{G1.} The codomain is declared as GR-shape geometry: a Lorentzian metric $g_{\mu\nu}$, a conformal class $[g]$, and the linearized Einstein equations $G_{\mu\nu} = 8\pi G\, T_{\mu\nu}$. The codomain category is \emph{differential geometry on a Lorentzian manifold}, an external category to the trace category $\mathcal{T}$.
\item \textbf{G2.} Four bank witnesses provide the transport. \texttt{check\_L\_HKM\_causal\_geometry} (causal order from $L_{\rm irr}$ uniquely determines $[g]$; HKM 1976). \texttt{check\_L\_metric\_from\_entanglement\_data} (full $g_{\mu\nu}$ from algebraic data, no spatial input). \texttt{check\_L\_spacetime\_emergence\_v2} (spacetime from $A_1$ with zero spatial postulates). \texttt{check\_L\_Einstein\_from\_entanglement} (linearized Einstein equations from $\delta Q = T\, dS$). All non-inverse: GR's structure is the output, not the input.
\item \textbf{G3.} The mathematical inputs are declared: manifold theory, the entanglement-entropy framework (von Neumann entropy + first law), causal-set order theory. No numerical PDG-style external constants enter at this layer; the inputs are structural-mathematical.
\item \textbf{G4 (structural).} The domain of validity is declared as the finite-physical regime ($C_\Gamma < \infty$) within causally connected regions. Excluded regimes are named: quantum-gravity scale (Planck regime, parked as RP5 long-term capstone), far-from-equilibrium spacetimes (named in Paper 25 H4 boundary), non-Lorentzian structural variants (out of scope by declaration).
\item \textbf{G5.} The Lorentzian metric is not consumed as input. The derivation produces $[g]$ from $L_{\rm irr}$ and $g_{\mu\nu}$ from algebraic data; GR's content is output, not input.
\item \textbf{G6 (structural).} The edge cases at the derivation's boundary are named. Each named edge case has an explicit parking destination: RP5 for quantum gravity, Paper 25 H4 for far-from-equilibrium, ``out of scope'' for non-Lorentzian variants. No edge case is left unnamed.
\end{itemize}

All six gates close in the structural-form interpretation. RP1 satisfies the export theorem under the structural-codomain extension.

\subsection{Status of the lemma extension}
\label{subsec:lemma-extension}

The route-status determination lemma extends to nine-for-nine across the codomain-transport program: the six articulated mass routes (W on-shell, bottom, top, charm, light-quark, charged-lepton) plus the three closed Reconstruction Programs (RP1 spacetime emergence, RP3 matter-antimatter asymmetry, RP6 inflation) all instantiate gate-availability patterns consistent with their named statuses.

Sufficiency of G1--G6 (no-G7-exists) and uniqueness of the status determination (no other status patterns possible) remain open as separate theorem programs. The empirical evidence is sufficient to commit to the lemma as a working principle; deductive proof is the next theorem layer.

\subsection{Schema extension}
\label{subsec:schema-extension}

The structural extension implies one focused modification to the codomain-transport schema (\texttt{apf/codomain\_transport\_schema.py}, banked at v9.6): a new per-route field \texttt{codomain\_type $\in$ \{numerical, structural\}}. The gate-pattern check runs the appropriate G4 / G6 interpretation per route based on this field. Existing routes are typed at landing: mass routes, RP3, RP6, and RP-CT.H0 are numerical; RP1 is structural; routes added in the future declare their type at registration.

This modification is a focused schema-extension task: one new field, one new check (\texttt{check\_T\_codomain\_type\_declared}), one expected bump. It is queued as a Tier-1 follow-on task and does not require any change to the v35 export-theorem statement itself.
